% Options for packages loaded elsewhere
\PassOptionsToPackage{unicode}{hyperref}
\PassOptionsToPackage{hyphens}{url}
\documentclass[
  a4paper,
]{ltjsarticle}
\usepackage{xcolor}
\usepackage[margin=25mm]{geometry}
\usepackage{amsmath,amssymb}
\setcounter{secnumdepth}{-\maxdimen} % remove section numbering
\usepackage{iftex}
\ifPDFTeX
  \usepackage[T1]{fontenc}
  \usepackage[utf8]{inputenc}
  \usepackage{textcomp} % provide euro and other symbols
\else % if luatex or xetex
  \usepackage{unicode-math} % this also loads fontspec
  \defaultfontfeatures{Scale=MatchLowercase}
  \defaultfontfeatures[\rmfamily]{Ligatures=TeX,Scale=1}
\fi
\usepackage{lmodern}
\ifPDFTeX\else
  % xetex/luatex font selection
\fi
% Use upquote if available, for straight quotes in verbatim environments
\IfFileExists{upquote.sty}{\usepackage{upquote}}{}
\IfFileExists{microtype.sty}{% use microtype if available
  \usepackage[]{microtype}
  \UseMicrotypeSet[protrusion]{basicmath} % disable protrusion for tt fonts
}{}
\makeatletter
\@ifundefined{KOMAClassName}{% if non-KOMA class
  \IfFileExists{parskip.sty}{%
    \usepackage{parskip}
  }{% else
    \setlength{\parindent}{0pt}
    \setlength{\parskip}{6pt plus 2pt minus 1pt}}
}{% if KOMA class
  \KOMAoptions{parskip=half}}
\makeatother
\usepackage{color}
\usepackage{fancyvrb}
\newcommand{\VerbBar}{|}
\newcommand{\VERB}{\Verb[commandchars=\\\{\}]}
\DefineVerbatimEnvironment{Highlighting}{Verbatim}{commandchars=\\\{\}}
% Add ',fontsize=\small' for more characters per line
\newenvironment{Shaded}{}{}
\newcommand{\AlertTok}[1]{\textcolor[rgb]{1.00,0.00,0.00}{\textbf{#1}}}
\newcommand{\AnnotationTok}[1]{\textcolor[rgb]{0.38,0.63,0.69}{\textbf{\textit{#1}}}}
\newcommand{\AttributeTok}[1]{\textcolor[rgb]{0.49,0.56,0.16}{#1}}
\newcommand{\BaseNTok}[1]{\textcolor[rgb]{0.25,0.63,0.44}{#1}}
\newcommand{\BuiltInTok}[1]{\textcolor[rgb]{0.00,0.50,0.00}{#1}}
\newcommand{\CharTok}[1]{\textcolor[rgb]{0.25,0.44,0.63}{#1}}
\newcommand{\CommentTok}[1]{\textcolor[rgb]{0.38,0.63,0.69}{\textit{#1}}}
\newcommand{\CommentVarTok}[1]{\textcolor[rgb]{0.38,0.63,0.69}{\textbf{\textit{#1}}}}
\newcommand{\ConstantTok}[1]{\textcolor[rgb]{0.53,0.00,0.00}{#1}}
\newcommand{\ControlFlowTok}[1]{\textcolor[rgb]{0.00,0.44,0.13}{\textbf{#1}}}
\newcommand{\DataTypeTok}[1]{\textcolor[rgb]{0.56,0.13,0.00}{#1}}
\newcommand{\DecValTok}[1]{\textcolor[rgb]{0.25,0.63,0.44}{#1}}
\newcommand{\DocumentationTok}[1]{\textcolor[rgb]{0.73,0.13,0.13}{\textit{#1}}}
\newcommand{\ErrorTok}[1]{\textcolor[rgb]{1.00,0.00,0.00}{\textbf{#1}}}
\newcommand{\ExtensionTok}[1]{#1}
\newcommand{\FloatTok}[1]{\textcolor[rgb]{0.25,0.63,0.44}{#1}}
\newcommand{\FunctionTok}[1]{\textcolor[rgb]{0.02,0.16,0.49}{#1}}
\newcommand{\ImportTok}[1]{\textcolor[rgb]{0.00,0.50,0.00}{\textbf{#1}}}
\newcommand{\InformationTok}[1]{\textcolor[rgb]{0.38,0.63,0.69}{\textbf{\textit{#1}}}}
\newcommand{\KeywordTok}[1]{\textcolor[rgb]{0.00,0.44,0.13}{\textbf{#1}}}
\newcommand{\NormalTok}[1]{#1}
\newcommand{\OperatorTok}[1]{\textcolor[rgb]{0.40,0.40,0.40}{#1}}
\newcommand{\OtherTok}[1]{\textcolor[rgb]{0.00,0.44,0.13}{#1}}
\newcommand{\PreprocessorTok}[1]{\textcolor[rgb]{0.74,0.48,0.00}{#1}}
\newcommand{\RegionMarkerTok}[1]{#1}
\newcommand{\SpecialCharTok}[1]{\textcolor[rgb]{0.25,0.44,0.63}{#1}}
\newcommand{\SpecialStringTok}[1]{\textcolor[rgb]{0.73,0.40,0.53}{#1}}
\newcommand{\StringTok}[1]{\textcolor[rgb]{0.25,0.44,0.63}{#1}}
\newcommand{\VariableTok}[1]{\textcolor[rgb]{0.10,0.09,0.49}{#1}}
\newcommand{\VerbatimStringTok}[1]{\textcolor[rgb]{0.25,0.44,0.63}{#1}}
\newcommand{\WarningTok}[1]{\textcolor[rgb]{0.38,0.63,0.69}{\textbf{\textit{#1}}}}
\usepackage{longtable,booktabs,array}
\newcounter{none} % for unnumbered tables
\usepackage{calc} % for calculating minipage widths
% Correct order of tables after \paragraph or \subparagraph
\usepackage{etoolbox}
\makeatletter
\patchcmd\longtable{\par}{\if@noskipsec\mbox{}\fi\par}{}{}
\makeatother
% Allow footnotes in longtable head/foot
\IfFileExists{footnotehyper.sty}{\usepackage{footnotehyper}}{\usepackage{footnote}}
\makesavenoteenv{longtable}
\setlength{\emergencystretch}{3em} % prevent overfull lines
\providecommand{\tightlist}{%
  \setlength{\itemsep}{0pt}\setlength{\parskip}{0pt}}
\usepackage{bookmark}
\IfFileExists{xurl.sty}{\usepackage{xurl}}{} % add URL line breaks if available
\urlstyle{same}
\hypersetup{
  pdftitle={第3章 複素平面読出し図 --- step0・終了時・凝縮中心拡大の3グリッド},
  pdfauthor={木原 範昭（WF System Co., Ltd.）　ORCID 0009-0004-6753-4020},
  hidelinks,
  pdfcreator={LaTeX via pandoc}}

\title{第3章 複素平面読出し図 ---
step0・終了時・凝縮中心拡大の3グリッド}
\author{木原 範昭（WF System Co., Ltd.）　ORCID 0009-0004-6753-4020}
\date{2026年9月5日　Version DOI 10.5281/zenodo.22317636}

\begin{document}
\maketitle

（N=3..40 段1+2+3 スイープ再現論文・第3章。総括論文と Concept DOI
10.5281/zenodo.22317635 を共有する。Version DOI
10.5281/zenodo.22317636。 式番号は第1章・第2章（式1〜式25）からの連番）

\subsection{1. 目的}\label{ux76eeux7684}

第2章のスイープが保存した状態 npz（Δτ=2π/N 走行）から、各辺の複素波 z\_e
∈ C を 複素平面に読み出した3枚のグリッド図------(1) step0、(2)
終了時（step500）、 (3)
終了時の最大角クラスター拡大------を生成する処理を記述する。3図の設計意図
（インフレーション図の始点・終点・終点内部を配置側から裏づける役割分担、および
「図の複素平面 ≠ 親平面 Π」の注意と橋渡しの恒等式）は\textbf{第2章
§2.6（式25）}で
与えられており、本章はその実装・実行・観察を固定する。データは読み出しのみで、
力学・保存データには一切触れない。

\subsection{2. 理論的背景}\label{ux7406ux8ad6ux7684ux80ccux666f}

状態と測定の数学は第2章（式22〜式25）に依る。本章の図固有の規約を式26〜式28
として 定義する。

\textbf{(式26) 重複計数（縮退の表記）} --- 描画対象の複素値集合 \{z\_e\}
を、実部・虚部の \textbf{小数第12位丸め}で同値類に分け、要素数 c
\textgreater{} 1 の類の位置に「xc」を注記する:

\begin{verbatim}
class(w) = ( round(Re w, 12), round(Im w, 12) )
\end{verbatim}

12桁丸めは倍精度の実効桁（〜16桁）より粗く、丸め誤差程度の差を同一視しつつ
物理的に有意な分離（本系列では相対 10\^{}-4
以上）は別扱いにする閾値である。
拡大図では丸めを第15位（式28）に締め、機械精度レベルの一致のみを「xn」と数える。

\textbf{(式27) パネルスケール規約} ---
各パネルの軸範囲は、そのパネルの全体振幅 r = max\_e
\textbar z\_e\textbar{} を用いて {[}−1.15r, +1.15r{]} の正方（equal
aspect）とし、
\textbf{目盛数値は実値のまま}表示する（正規化した値に貼り替えない）。原点から各点への
線分は波の偏角と振幅を同時に読むための補助線である。

\textbf{(式28) 最大角クラスター抽出（拡大図の算法）} ---
終了時状態を全体振幅で規格化した 座標で\textbf{粗解像度 1/100}
に丸めて群化し、最大要素数の群を拡大対象とする:

\begin{verbatim}
key(w) = ( round(Re w / amp, 2), round(Im w / amp, 2) ),   amp = max_e |z_e|
C* = argmax_{key} |{ w : key(w) = key }|
中心 c = mean(C*),  割れ幅 spread = max_{w∈C*} |w − c|,  窓幅 = 1.4 × spread
\end{verbatim}

パネル題に本数 \textbar C*\textbar・spread・amp を印字する。spread
は「クラスターが厳密な一点
（厳密縮退）か有限幅の束か」を定量化する（第2章 §2.6
の表・拡大図の行）。

\subsection{3. 実装方法}\label{ux5b9fux88c5ux65b9ux6cd5}

\begin{itemize}
\tightlist
\item
  図化プログラム
  \texttt{plot\_complex\_plane\_N3\_N40\_stage123\_v1.py}（本パッケージ同梱、
  読み出しのみ）。描画様式は本シリーズの既存図化プログラム
  （\texttt{complex\_plane\_readout\_step0\_step2000\_20260904/plot\_complex\_plane\_step0\_step2000.py}
  のグリッド様式、および
  \texttt{自発的分裂予備実験\_v1/N40\_state\_readout\_20260904/\ plot\_complex\_plane\_N40\_v1.py}
  のクラスター拡大算法）を N=3..40 の 8×5 グリッドに
  拡張したもので、算法自体は同一である（様式の系譜を変えない）。
\item
  入力は第2章の成果物
  \texttt{results/hm\_N\{N\}\_den\_\{N\}\_states\_500.npz}
  のみ。\textbf{den=N （Δτ=2π/N）の走行を各 N の代表として用いる}（step0
  は分母に依らず同一なので
  代表選択は終了時にのみ効く。他の分母の終了時配置は保存済み npz
  から同様に 読み出せる）。
\item
  出力は PNG 3枚。データの書き換え・再生成はない。
\end{itemize}

\subsection{4. 詳細設計}\label{ux8a73ux7d30ux8a2dux8a08}

\subsubsection{4.1 全体フロー}\label{ux5168ux4f53ux30d5ux30edux30fc}

\begin{verbatim}
[初期化]   入力フォルダの決定（results/）
[ループ処理] draw_grid(step=0)  : N=3..40 の各パネルに step0 状態を描画（式26・式27）
           draw_grid(step=500): 同、終了時状態
           拡大グリッド      : N=3..40 の各パネルに最大角クラスター（式28）
[終了処理] 各図を savefig、未使用2パネルを off
\end{verbatim}

\subsubsection{4.2
全体データフロー}\label{ux5168ux4f53ux30c7ux30fcux30bfux30d5ux30edux30fc}

\begin{itemize}
\tightlist
\item
  \textbf{入力}:
  \texttt{results/hm\_N\{N\}\_den\_\{N\}\_states\_500.npz} ×
  38（第2章成果物・読み出しのみ）。 使用フィールド:
  \texttt{Z}（501×M）、\texttt{denominator}、\texttt{steps}（整合性
  assert 用）
\item
  \textbf{パラメータ}（すべてプログラム内定数）:

  \begin{itemize}
  \tightlist
  \item
    丸め桁: 12（式26）／15（拡大図内、式28）／粗解像度 2桁（式28の
    1/100）
  \item
    スケール係数 1.15（式27）、窓係数 1.4（式28）
  \item
    グリッド 8×5（38使用・2 off）、dpi=180
  \end{itemize}
\item
  \textbf{出力}:

  \begin{itemize}
  \tightlist
  \item
    \texttt{fig\_complex\_plane\_step0\_N3\_N40\_stage123.png}
  \item
    \texttt{fig\_complex\_plane\_final\_N3\_N40\_stage123.png}
  \item
    \texttt{fig\_complex\_plane\_final\_zoom\_N3\_N40\_stage123.png}
  \end{itemize}
\end{itemize}

\subsubsection{4.3 個別処理}\label{ux500bux5225ux51e6ux7406}

\paragraph{4.3.1
初期化（読み込みと整合性検査）}\label{ux521dux671fux5316ux8aadux307fux8fbcux307fux3068ux6574ux5408ux6027ux691cux67fb}

\begin{Shaded}
\begin{Highlighting}[]
    \DecValTok{18}  \KeywordTok{def}\NormalTok{ load(N, step):}
    \DecValTok{19}\NormalTok{      d }\OperatorTok{=}\NormalTok{ np.load(os.path.join(IN, }\SpecialStringTok{f\textquotesingle{}hm\_N}\SpecialCharTok{\{}\NormalTok{N}\SpecialCharTok{\}}\SpecialStringTok{\_den\_}\SpecialCharTok{\{}\NormalTok{N}\SpecialCharTok{\}}\SpecialStringTok{\_states\_500.npz\textquotesingle{}}\NormalTok{))}
    \DecValTok{20}      \ControlFlowTok{assert} \BuiltInTok{int}\NormalTok{(d[}\StringTok{\textquotesingle{}denominator\textquotesingle{}}\NormalTok{]) }\OperatorTok{==}\NormalTok{ N }\KeywordTok{and} \BuiltInTok{int}\NormalTok{(d[}\StringTok{\textquotesingle{}steps\textquotesingle{}}\NormalTok{]) }\OperatorTok{==} \DecValTok{500}
    \DecValTok{21}      \ControlFlowTok{return}\NormalTok{ np.asarray(d[}\StringTok{\textquotesingle{}Z\textquotesingle{}}\NormalTok{][step], dtype}\OperatorTok{=}\NormalTok{np.complex128)}
\end{Highlighting}
\end{Shaded}

\begin{itemize}
\tightlist
\item
  20行の assert は「意図したファイル（den=N・500
  step）を読んでいる」ことの検査で、 不一致なら AssertionError
  で停止する（描画は行われない）。
\end{itemize}

\paragraph{4.3.2 ループ処理（1）:
グリッド描画（式26・式27）}\label{ux30ebux30fcux30d7ux51e6ux74061-ux30b0ux30eaux30c3ux30c9ux63cfux753bux5f0f26ux5f0f27}

\begin{Shaded}
\begin{Highlighting}[]
    \DecValTok{26}      \ControlFlowTok{for}\NormalTok{ k, N }\KeywordTok{in} \BuiltInTok{enumerate}\NormalTok{(}\BuiltInTok{range}\NormalTok{(}\DecValTok{3}\NormalTok{, }\DecValTok{41}\NormalTok{)):}
    \DecValTok{27}\NormalTok{          ax }\OperatorTok{=}\NormalTok{ axs[k]}
    \DecValTok{28}\NormalTok{          z }\OperatorTok{=}\NormalTok{ load(N, step)}
    \DecValTok{29}\NormalTok{          M }\OperatorTok{=}\NormalTok{ N }\OperatorTok{*}\NormalTok{ (N }\OperatorTok{{-}} \DecValTok{1}\NormalTok{) }\OperatorTok{//} \DecValTok{2}
    \DecValTok{30}          \ControlFlowTok{assert}\NormalTok{ z.size }\OperatorTok{==}\NormalTok{ M}
    \DecValTok{31}          \ControlFlowTok{for}\NormalTok{ w }\KeywordTok{in}\NormalTok{ z:}
    \DecValTok{32}\NormalTok{              ax.plot([}\FloatTok{0.0}\NormalTok{, w.real], [}\FloatTok{0.0}\NormalTok{, w.imag], color}\OperatorTok{=}\StringTok{\textquotesingle{}tab:blue\textquotesingle{}}\NormalTok{, linewidth}\OperatorTok{=}\FloatTok{0.7}\NormalTok{, alpha}\OperatorTok{=}\FloatTok{0.6}\NormalTok{)}
    \DecValTok{33}\NormalTok{          ax.plot(z.real, z.imag, }\StringTok{\textquotesingle{}o\textquotesingle{}}\NormalTok{, ms}\OperatorTok{=}\FloatTok{2.5}\NormalTok{, color}\OperatorTok{=}\StringTok{\textquotesingle{}tab:red\textquotesingle{}}\NormalTok{, alpha}\OperatorTok{=}\FloatTok{0.85}\NormalTok{, linestyle}\OperatorTok{=}\StringTok{\textquotesingle{}none\textquotesingle{}}\NormalTok{)}
    \DecValTok{34}\NormalTok{          cnt }\OperatorTok{=}\NormalTok{ Counter((}\BuiltInTok{round}\NormalTok{(}\BuiltInTok{float}\NormalTok{(w.real), }\DecValTok{12}\NormalTok{), }\BuiltInTok{round}\NormalTok{(}\BuiltInTok{float}\NormalTok{(w.imag), }\DecValTok{12}\NormalTok{)) }\ControlFlowTok{for}\NormalTok{ w }\KeywordTok{in}\NormalTok{ z)}
    \DecValTok{35}          \ControlFlowTok{for}\NormalTok{ (a, b), c }\KeywordTok{in}\NormalTok{ cnt.items():}
    \DecValTok{36}              \ControlFlowTok{if}\NormalTok{ c }\OperatorTok{\textgreater{}} \DecValTok{1}\NormalTok{:}
    \DecValTok{37}\NormalTok{                  ax.annotate(}\SpecialStringTok{f\textquotesingle{}x}\SpecialCharTok{\{}\NormalTok{c}\SpecialCharTok{\}}\SpecialStringTok{\textquotesingle{}}\NormalTok{, (a, b), textcoords}\OperatorTok{=}\StringTok{\textquotesingle{}offset points\textquotesingle{}}\NormalTok{, xytext}\OperatorTok{=}\NormalTok{(}\DecValTok{3}\NormalTok{, }\DecValTok{3}\NormalTok{),}
    \DecValTok{38}\NormalTok{                              fontsize}\OperatorTok{=}\DecValTok{5}\NormalTok{, color}\OperatorTok{=}\StringTok{\textquotesingle{}black\textquotesingle{}}\NormalTok{)}
    \DecValTok{39}\NormalTok{          r }\OperatorTok{=} \BuiltInTok{float}\NormalTok{(np.}\BuiltInTok{abs}\NormalTok{(z).}\BuiltInTok{max}\NormalTok{())}
    \DecValTok{40}\NormalTok{          lim }\OperatorTok{=}\NormalTok{ r }\OperatorTok{*} \FloatTok{1.15} \ControlFlowTok{if}\NormalTok{ r }\OperatorTok{\textgreater{}} \DecValTok{0} \ControlFlowTok{else} \FloatTok{1.0}
    \DecValTok{41}\NormalTok{          ax.set\_xlim(}\OperatorTok{{-}}\NormalTok{lim, lim)}\OperatorTok{;}\NormalTok{ ax.set\_ylim(}\OperatorTok{{-}}\NormalTok{lim, lim)}\OperatorTok{;}\NormalTok{ ax.set\_aspect(}\StringTok{\textquotesingle{}equal\textquotesingle{}}\NormalTok{)}
\end{Highlighting}
\end{Shaded}

\begin{itemize}
\tightlist
\item
  31-32行: 原点からの線分、33行: 点、34-38行: 式26の重複注記、39-41行:
  式27の スケール規約（目盛は実値のまま。46行の
  \texttt{ticklabel\_format} は表記の指数化のみ）。
\end{itemize}

\paragraph{4.3.3 ループ処理（2）:
最大角クラスター拡大（式28）}\label{ux30ebux30fcux30d7ux51e6ux74062-ux6700ux5927ux89d2ux30afux30e9ux30b9ux30bfux30fcux62e1ux5927ux5f0f28}

\begin{Shaded}
\begin{Highlighting}[]
    \DecValTok{67}\NormalTok{      z }\OperatorTok{=}\NormalTok{ load(N, }\DecValTok{500}\NormalTok{)}
    \DecValTok{68}\NormalTok{      amp }\OperatorTok{=} \BuiltInTok{float}\NormalTok{(np.}\BuiltInTok{abs}\NormalTok{(z).}\BuiltInTok{max}\NormalTok{())}
    \DecValTok{69}\NormalTok{      coarse }\OperatorTok{=}\NormalTok{ \{\}}
    \DecValTok{70}      \ControlFlowTok{for}\NormalTok{ w }\KeywordTok{in}\NormalTok{ z:}
    \DecValTok{71}\NormalTok{          key }\OperatorTok{=}\NormalTok{ (}\BuiltInTok{round}\NormalTok{(}\BuiltInTok{float}\NormalTok{(w.real) }\OperatorTok{/}\NormalTok{ amp, }\DecValTok{2}\NormalTok{), }\BuiltInTok{round}\NormalTok{(}\BuiltInTok{float}\NormalTok{(w.imag) }\OperatorTok{/}\NormalTok{ amp, }\DecValTok{2}\NormalTok{))}
    \DecValTok{72}\NormalTok{          coarse.setdefault(key, []).append(w)}
    \DecValTok{73}\NormalTok{      mem }\OperatorTok{=} \BuiltInTok{max}\NormalTok{(coarse.values(), key}\OperatorTok{=}\BuiltInTok{len}\NormalTok{)}
    \DecValTok{74}\NormalTok{      zz }\OperatorTok{=}\NormalTok{ np.array(mem)}
    \DecValTok{75}\NormalTok{      c }\OperatorTok{=}\NormalTok{ zz.mean()}
    \DecValTok{76}\NormalTok{      dev }\OperatorTok{=}\NormalTok{ np.}\BuiltInTok{abs}\NormalTok{(zz }\OperatorTok{{-}}\NormalTok{ c)}
    \DecValTok{77}\NormalTok{      spread }\OperatorTok{=} \BuiltInTok{float}\NormalTok{(dev.}\BuiltInTok{max}\NormalTok{())}
    \DecValTok{78}\NormalTok{      win }\OperatorTok{=}\NormalTok{ spread }\OperatorTok{*} \FloatTok{1.4} \ControlFlowTok{if}\NormalTok{ spread }\OperatorTok{\textgreater{}} \DecValTok{0} \ControlFlowTok{else}\NormalTok{ amp }\OperatorTok{*} \FloatTok{1e{-}12}
\NormalTok{...}
    \DecValTok{90}\NormalTok{      ax.set\_title(}\SpecialStringTok{f\textquotesingle{}N=}\SpecialCharTok{\{}\NormalTok{N}\SpecialCharTok{\}}\SpecialStringTok{: }\SpecialCharTok{\{}\BuiltInTok{len}\NormalTok{(zz)}\SpecialCharTok{\}}\SpecialStringTok{ waves, dev=}\SpecialCharTok{\{}\NormalTok{spread}\SpecialCharTok{:.2e\}}\SpecialStringTok{ (|z|max=}\SpecialCharTok{\{}\NormalTok{amp}\SpecialCharTok{:.2e\}}\SpecialStringTok{)\textquotesingle{}}\NormalTok{, fontsize}\OperatorTok{=}\DecValTok{7}\NormalTok{)}
\end{Highlighting}
\end{Shaded}

\begin{itemize}
\tightlist
\item
  70-73行: 式28の群化と最大クラスター選択、75-78行: 中心・割れ幅・窓幅、
  90行: パネル題への本数・spread・amp の印字（§6
  の観察値の出所はこの印字である）。
\end{itemize}

\paragraph{4.3.4
終了処理・例外的挙動（数式化しない要素）}\label{ux7d42ux4e86ux51e6ux7406ux4f8bux5916ux7684ux6319ux52d5ux6570ux5f0fux5316ux3057ux306aux3044ux8981ux7d20}

\begin{itemize}
\tightlist
\item
  40行 \texttt{lim\ =\ r*1.15\ if\ r\ \textgreater{}\ 0\ else\ 1.0}:
  全点が原点の退化ケースへのフォールバック （本データでは発生しない）。
\item
  78行
  \texttt{win\ =\ spread*1.4\ if\ spread\ \textgreater{}\ 0\ else\ amp*1e-12}:
  クラスターが1点のみ （spread=0）の場合の窓幅フォールバック（N=3
  で発生: 最大クラスターが1波）。
\item
  50-51行・91-92行: 8×5=40 パネル中、未使用の 39・40 番目を
  \texttt{axis(\textquotesingle{}off\textquotesingle{})}。
\item
  例外処理は assert（20・30行）のみで、それ以外の分岐はない。
\end{itemize}

\subsection{5. 実行結果}\label{ux5b9fux884cux7d50ux679c}

\subsubsection{5.1
再現コマンド}\label{ux518dux73feux30b3ux30deux30f3ux30c9}

\begin{Shaded}
\begin{Highlighting}[]
\BuiltInTok{cd}\NormalTok{ N3\_N40\_stage123\_sweep\_20260905}
\ExtensionTok{python3}\NormalTok{ plot\_complex\_plane\_N3\_N40\_stage123\_v1.py}
\CommentTok{\# または ./run\_all.sh の最終段として実行される}
\end{Highlighting}
\end{Shaded}

\subsubsection{5.2 実行環境}\label{ux5b9fux884cux74b0ux5883}

\begin{itemize}
\tightlist
\item
  Python 3.9.6（\texttt{.venv/bin/python3}）、numpy
  2.0.2、matplotlib（Agg 描画）
\item
  macOS 26.3.1（arm64）
\end{itemize}

\subsubsection{5.3 実行時間}\label{ux5b9fux884cux6642ux9593}

3図の生成合計で\textbf{約1分以内}（支配項は 38×2 回の npz 読込と 780
本×38 パネルの線分描画）。

\subsubsection{5.4 検証ゲート}\label{ux691cux8a3cux30b2ux30fcux30c8}

{\def\LTcaptype{none} % do not increment counter
\begin{longtable}[]{@{}
  >{\raggedright\arraybackslash}p{(\linewidth - 6\tabcolsep) * \real{0.2500}}
  >{\raggedright\arraybackslash}p{(\linewidth - 6\tabcolsep) * \real{0.2500}}
  >{\raggedright\arraybackslash}p{(\linewidth - 6\tabcolsep) * \real{0.2500}}
  >{\raggedright\arraybackslash}p{(\linewidth - 6\tabcolsep) * \real{0.2500}}@{}}
\toprule\noalign{}
\begin{minipage}[b]{\linewidth}\raggedright
ゲート
\end{minipage} & \begin{minipage}[b]{\linewidth}\raggedright
合格条件
\end{minipage} & \begin{minipage}[b]{\linewidth}\raggedright
実測
\end{minipage} & \begin{minipage}[b]{\linewidth}\raggedright
合否
\end{minipage} \\
\midrule\noalign{}
\endhead
\bottomrule\noalign{}
\endlastfoot
G1: 入力整合 & 各読込で denominator==N・steps==500・z.size==M の assert
通過 & 全38N×3図で通過（AssertionError なし） & \textbf{PASS} \\
G2: 完走 & \texttt{ALL\ DONE} 出力・PNG 3枚生成 & 確認 &
\textbf{PASS} \\
\end{longtable}
}

（入力 npz 自体の bit 系譜は第2章 G1〔全228走行の Z{[}0{]} が静的親と
bit 一致〕で担保済み）

\subsubsection{5.5 データ}\label{ux30c7ux30fcux30bf}

{\def\LTcaptype{none} % do not increment counter
\begin{longtable}[]{@{}
  >{\raggedright\arraybackslash}p{(\linewidth - 2\tabcolsep) * \real{0.5000}}
  >{\raggedright\arraybackslash}p{(\linewidth - 2\tabcolsep) * \real{0.5000}}@{}}
\toprule\noalign{}
\begin{minipage}[b]{\linewidth}\raggedright
項目
\end{minipage} & \begin{minipage}[b]{\linewidth}\raggedright
内容
\end{minipage} \\
\midrule\noalign{}
\endhead
\bottomrule\noalign{}
\endlastfoot
入力 & \texttt{results/hm\_N\{N\}\_den\_\{N\}\_states\_500.npz} ×
38（第2章成果物、読み出しのみ） \\
出力先 & パッケージ直下 \\
step0 図 &
\texttt{fig\_complex\_plane\_step0\_N3\_N40\_stage123.png}（632,949
bytes） \\
終了時図 &
\texttt{fig\_complex\_plane\_final\_N3\_N40\_stage123.png}（4,582,367
bytes） \\
拡大図 &
\texttt{fig\_complex\_plane\_final\_zoom\_N3\_N40\_stage123.png}（494,906
bytes） \\
SHA256 & 同梱 \texttt{SHA256SUMS.txt} を正本とする \\
\end{longtable}
}

\subsubsection{5.6 図化}\label{ux56f3ux5316}

上記3枚（8×5 グリッド、各パネル equal aspect・実値目盛）。図の読み方と
インフレーション図との対応は第2章 §2.6 の表の通り。

\subsection{6.
実行分析（客観的報告と観察のみ。数値の出所は図中印字＝本プログラムの出力）}\label{ux5b9fux884cux5206ux6790ux5ba2ux89b3ux7684ux5831ux544aux3068ux89b3ux5bdfux306eux307fux6570ux5024ux306eux51faux6240ux306fux56f3ux4e2dux5370ux5b57ux672cux30d7ux30edux30b0ux30e9ux30e0ux306eux51faux529b}

\begin{enumerate}
\def\labelenumi{\arabic{enumi}.}
\tightlist
\item
  \textbf{step0 図}: 全 N=3..40
  のパネルで、対蹠2ペア（4束）の星型が現れる。束は動径方向に
  振幅の広がりを持つ。式26の重複注記は小 N の一部（N=4 の x4、N=5〜9 の
  x2〜x3 など） に現れ、N が大きくなると厳密重複は消える。
\item
  \textbf{終了時図}:
  全パネルで星型は消失し、ほぼ等振幅のリング状配置（小 N では疎な
  スポーク）になる。式26（12桁丸め）の重複注記は小 N の少数（N=4, 5 の
  x2）を除き 現れない。
\item
  \textbf{拡大図}: 最大角クラスターの本数は 1〜10（N=3: 1波、N=39, 40:
  10波）、割れ幅 spread は N≥6 で 10\textsuperscript{-7〜10}-4
  のオーダー（例: N=40 で
  dev=2.04e-04、\textbar z\textbar max=3.64e-02、 相対
  \textasciitilde5.6×10\^{}-3）。N=4, 5 のクラスター（2波）のみ dev が
  10\textsuperscript{-16〜10}-10 と機械精度 レベルで一致する。
\item
  観察のまとめ（事実のみ）: 終了時の角クラスターは、N≥6
  の全域で「厳密な同一複素数 への凝縮」ではなく有限幅の束である。第2章
  §2.6 の役割分担の通り、これは H⊥/H が
  測らない終状態の微細組織の記録である。
\end{enumerate}

\begin{center}\rule{0.5\linewidth}{0.5pt}\end{center}

（第3章おわり。総括論文は別途）

\end{document}
